\section{Introduction}
\label{sec:intro}
%
A central challenge in autonomous machine-learning research is to translate experimental feedback into effective model improvements and accumulate useful knowledge across successive experiments. Research agents tackle this challenge with code, tools, and feedback for experiments, language-model post-training, and training-data research~\citep{mlagentbench,mlebench,rebench,aiscientist,posttrainbench,rsibenchdata}. Beyond text-only post-training, agents must connect evidence in source media with language and other modalities. Clutter, background sounds, and incomplete descriptions in media collected in the wild make relevant evidence difficult to isolate. This raises a broader question: \emph{Can research agents autonomously improve how models see, hear, and reason about multimodal inputs---and sustain those improvements as experimentation continues?}

%
Autonomous multimodal post-training couples three challenges. \textbf{(1) Multimodal Error Diagnosis.} In visual reasoning, chain-of-thought prompting can degrade performance on perception-heavy tasks~\citep{mmecot}. Agents must inspect source media to distinguish missed or misinterpreted evidence from faulty inference before choosing an intervention. \textbf{(2) Data Curation and Modality Alignment.} Selecting visual instructions for vulnerable sample groups can improve robustness~\citep{ards}. Agents must also preserve semantic and temporal alignment: labels should match observable content, and paired audio-video segments should retain their timing. \textbf{(3) Budgeted Multimodal Experimentation.} Autonomous language-model post-training evaluations report early termination without valid final checkpoints~\citep{posttrainbench}. Multimodal runs must also budget for frame sampling and audio processing alongside training, candidate evaluation, and final delivery. Effective research must therefore connect media evidence and alignment decisions to the improvements ultimately delivered.

%
We introduce \benchmark{} (Multimodal Post-Training Benchmark)\footnote{Benchmark code and evaluation setup: \url{https://github.com/sod1010/MMPostTrainBench}.}, which evaluates autonomous research on eight tasks spanning image, audio, video, and joint audio-visual understanding, as well as multimodal code tasks (Table~\ref{tab:tasks}; Figure~\ref{fig:protocol}). Task queries specify a target capability and benchmark, such as audio reasoning on MMAR, and require training pipelines to process the corresponding media. Under a common resource budget, agents select data sources, modality mixtures, supervision, and training strategies in response to development feedback. Harbor-compatible task packages provide this feedback during research and reserve separate held-out tests for final evaluation. Target and non-target evaluations assess whether a targeted update preserves other capabilities.

%
We assess \textbf{model outcomes}, \textbf{iterative model improvement}, and \textbf{research integrity}. Across six research-agent configurations on eight tasks, 52.1\% of model--task means fall below the base, and submissions also exhibit non-target regressions. Figure~\ref{fig:key-results} summarizes final outcomes and best-so-far candidate improvement. Final submissions trail the best evaluated candidates by up to 5.38 percentage points across recorded candidate pools. Integrity checks flag suspected misuse of evaluation data during training. Producing a strong candidate, selecting it reliably, and retaining other multimodal capabilities thus remain distinct requirements.

%
These findings motivate our proposed multimodal research framework, \harness{}. Building on persistent research orchestration~\citep{arbor}, it augments existing code-agent runtimes with \textbf{multimodal evidence reasoning}, \textbf{hierarchical memory}, and \textbf{evaluation-guided model selection}. Observations are anchored to source media and linked to hypotheses, interventions, and development outcomes; memory carries successful recipes and failure records across rounds; model selection retains a candidate for final delivery. These mechanisms connect diagnosis, accumulated knowledge, and submission decisions within a seven-step research loop. In the evaluated configurations, adding \harness{} improves submitted-model accuracy by up to 7.75 percentage points for Claude Opus 4.8 with Claude Code and 2.33 points for GPT-5.6-sol with Codex.

Our contributions are threefold: \textbf{(1)} We introduce \benchmark{}, which evaluates autonomous post-training across eight multimodal tasks through model outcomes, iterative model improvement, and research integrity. \textbf{(2)} We characterize non-target regressions, candidate degradation, and submission-selection failures in current research agents. \textbf{(3)} We propose \harness{}, which combines multimodal evidence reasoning, hierarchical memory, and evaluation-guided model selection to improve submitted models within existing code-agent runtimes.

%
%
%
%
%
%
%
%
%
%
%
%
%
%
%
